\subsection{Exact generic ranks on every principally polarized abelian surface in characteristic two}
\label{sec:DE-A3-ppav-ranks}

The next result gives a structural formula for this entire class of
surfaces.  In particular the ordinary elliptic factor of a mixed product
may have an arbitrary algebraic $j$-invariant.  The proof does not make
an assertion about a general smooth projective ambient variety.

\begin{lemma}[The four-parameter alternating theta tensor]
\label{lem:DE-A3-theta-pfaffian}
\status{Proved; root and monodromy full independent audits passed.}
Work over an algebraically closed field of characteristic two.  Index a
nine-dimensional space $W$ by the points of $I=\mathbf F_3^2$, and let
$\mathcal D=\mathbf P^1(\mathbf F_3)$ be its four line directions.
For $i\in\mathcal D$ put
\[
 P_i(s)=\sum_{\substack{\ell\subset I\text{ affine line}\\
                         \operatorname{dir}(\ell)=i}}
                           \prod_{a\in\ell}s_a,
 \qquad P_*(s)=\sum_{a\in I}s_a^3.
\]
Let $T$ be the alternating trilinear form whose value on three distinct
basis vectors is $\alpha_i$ when their indices form a line of direction
$i$, and zero otherwise.  Let $M(s)$ be its contraction matrix:
\[
             M(s)_{ab}=T(s,e_b,e_a).
\]
Then $M(s)s=0$, and the vector of its eight-by-eight principal
Pfaffians is $q_\alpha(s)s$, where
\begin{equation}\label{eq:DE-A3-pfaffian-cubic}
 q_\alpha(s)=
   \left(\prod_{i\in\mathcal D}\alpha_i\right)P_*(s)
   +\sum_{i\in\mathcal D}
       \alpha_i\left(\sum_{j\ne i}\alpha_j^3\right)P_i(s).
\end{equation}
Consequently
\[
                \operatorname{adj}M(s)=q_\alpha(s)^2ss^t.
\]
For a nonzero coefficient vector $\alpha$, the cubic $q_\alpha$ vanishes
identically precisely in either of the following cases: exactly one
coefficient is nonzero; or exactly three are nonzero and their cubes are
equal.  These are forty points in the projective coefficient space.
In each case $W$ is a direct sum of three three-dimensional spaces on
which $T$ is the sum of three nonzero volume forms.  In particular it
has a nontrivial projector $D$ with
\[
                         T(Ds,t,u)=T(s,Dt,u).
\]
\end{lemma}
\begin{proof}
All identities may first be checked over the polynomial ring in the
$\alpha_i$ and $s_a$, and then specialized.  Alternation gives $M(s)s=0$.
The Pfaffian kernel identity says that its vector $v$ of principal
Pfaffians also lies in the kernel.  The universal matrix has rank eight:
at $s=e_0$, deleting the zero index leaves four two-by-two alternating
blocks with entries $\alpha_i$, and their Pfaffian is
$\prod_i\alpha_i$.  Thus over the fraction field $v=q s$.
Since the coordinates $s_a$ are pairwise coprime, the identities
$v_as_b=v_bs_a$ imply that each $s_a$ divides $v_a$ in the polynomial
ring.  Hence $q$ is a polynomial, cubic in $s$ and quartic in $\alpha$.

The finite level-three Heisenberg group preserves $T$.  Its generators
are translations of the index set, the diagonal modulations
$e_a\mapsto\zeta^{y\cdot a}e_a$, and scalar multiplication by
$\zeta$, where $\zeta^3=1$ and $\zeta\ne1$.  Each generator has
determinant one.  The transformation rule for principal Pfaffians of
an odd alternating matrix, or the identification
$\bigwedge^8W^*\simeq W\otimes\det(W)^*$, shows that $v$ transforms
as $s$ does.  Therefore $q$ is invariant under this group.
Modulation invariance permits only cubic monomials whose three indices,
with repetition, sum to zero.  Such a monomial is either $s_a^3$ or the
product of the three points on an affine line: if two indices coincide,
the third must coincide as well.  Translation invariance then shows
that the five polynomials $P_*,P_i$ are a basis for these invariant
cubics.

At $s=e_0$ the preceding four-block calculation gives the coefficient
$\prod_i\alpha_i$ of $P_*$.  To determine a line coefficient take the
vertical line $x=0$ and specialize its three section coordinates to
$x_0,x_1,x_2$, setting the other six to zero.  Deleting the index
$(0,0)$ leaves its two vertical companions as one block, with entry
$\alpha_\infty x_0$, and a six-by-six block with Pfaffian $\det B$, where
\[
 B=\begin{pmatrix}
 \alpha_0x_0&\alpha_1x_2&\alpha_2x_1\\
 \alpha_2x_2&\alpha_0x_1&\alpha_1x_0\\
 \alpha_1x_1&\alpha_2x_0&\alpha_0x_2
 \end{pmatrix}.
\]
Its six-term determinant is
\[
 \det B=\alpha_0\alpha_1\alpha_2(x_0^3+x_1^3+x_2^3)
       +(\alpha_0^3+\alpha_1^3+\alpha_2^3)x_0x_1x_2.
\]
Dividing the principal Pfaffian by $x_0$ proves the coefficient of
$P_\infty$.  Linear changes of the index coordinates permute the four
directions, proving every coefficient in
\eqref{eq:DE-A3-pfaffian-cubic}.  The adjugate formula is the usual
odd alternating identity $\operatorname{adj}M=vv^t$, also a polynomial
identity, so it holds at every specialization.

The five orbit sums have disjoint monomial supports.  Thus $q=0$ is
equivalent to
\[
 \prod_i\alpha_i=0,\qquad
 \alpha_i\sum_{j\ne i}\alpha_j^3=0\quad(i\in\mathcal D).
\]
If $r$ coefficients are nonzero, the last equations say that all their
cubes equal the same nonzero number $u$ and that $ru=u$ in the
characteristic-two field.  Therefore $r$ is odd, and the product equation
leaves $r=1$ or $r=3$.  Projectively the first type gives four points;
the second gives four choices of the missing direction and $3^2$ choices
of the two ratios in $\mu_3$, hence thirty-six points.

For $r=1$ the three parallel affine lines partition the basis into
three orthogonal three-spaces, and $T$ is a nonzero volume form on each.
For $r=3$, scale $T$ so that its three nonzero coefficients belong to
$\mu_3$.  A diagonal basis change
\[
                  e_a\longmapsto\zeta^{Q(a)}e_a
\]
for a homogeneous quadratic form $Q$ on $\mathbf F_3^2$ multiplies the
coefficient in direction $v$ by $\zeta^{-Q(v)}$.  Indeed
$\sum_{t\in\mathbf F_3}Q(a+tv)=-Q(v)$.
Evaluations of binary quadratics on any three distinct projective
directions are independent: a nonzero binary quadratic cannot vanish at
three projective points.  One can therefore make the three coefficients
equal to one.  A linear permutation of the index set makes the missing
direction vertical.

Here is the remaining change of basis explicitly.  Start from the
one-direction form that is the sum of volume forms on the three
horizontal lines.  Apply Fourier change of basis in the second index:
\[
              f_{x,u}=\sum_{y\in\mathbf F_3}\zeta^{uy}e_{x,y}.
\]
This matrix is invertible; its inverse has entries $\zeta^{-uy}$,
because $3=1$ in the coefficient field.  Its transformed tensor is
nonzero exactly when the three first indices are distinct and the
three second indices sum to zero.  The coefficient is
$\sum_y\zeta^{y(u_1+u_2+u_3)}$, which equals one or zero accordingly.
These triples are exactly the affine lines in the three nonvertical
directions.  Thus the normalized three-direction form is equivalent
to the one-direction form.  Its three orthogonal projectors transport
back to those asserted in the lemma.  The displayed self-adjointness
identity is immediate from orthogonality of the three summands.
\end{proof}

\begin{proposition}[Every degree-one relative Frobenius is generically bijective]
\label{prop:DE-A3-ppav-degree-one}
\status{Proved; root and monodromy full independent audits passed.}
Let $(A,\Theta)$ be any principally polarized abelian surface in
characteristic two and put $L=\OO_A(3\Theta)$.  The relative coherent
Frobenius for the geometric generic smooth member of $|L|$ is bijective
on its eight-dimensional space.  If $f_A$ is the $2$-rank of $A$, that
curve has genus ten and $2$-rank $8+f_A$.
\end{proposition}
\begin{proof}
Work over an algebraic closure and choose a symmetric theta
representative.  Any other representative is a translate, giving an
isomorphic geometric family.  The typed theta construction of
Theorem~\ref{thm:ppav-level3-Cartier-symmetric} applies to every $A$,
with relative Cartier pairing
\[
 \beta:W\otimes W\longrightarrow W^{(2)},\qquad
 W=H^0(A,L),\quad\dim W=9.
\]
Its theta basis identifies the coefficient-square representation with
the contragredient of the finite Heisenberg representation; the resulting
trilinear tensor $T$ is fully symmetric and has diagonal coefficient
given by Cartier on an invariant top form.  All Frobenius twists and
the restriction to the finite central $\mu_3$ are as in that theorem.

If $A$ is ordinary, the independent symmetric diagonal gives a nonzero
full Hasse determinant by that theorem.  Its relative rank-eight
subbundle is consequently generically bijective as well.  We now
consider the other case.  Cartier on the invariant two-form is zero,
so $T(e_a,e_a,e_b)=0$ for every pair of basis indices.  Thus $T$ is
alternating.  Modulation invariance and index translations put it
exactly in the four-parameter form of
Lemma~\ref{lem:DE-A3-theta-pfaffian}.

We show that its cubic $q_\alpha$ cannot vanish.  Relative Frobenius
$F:A\to A^{(2)}$ is finite flat of degree four, and
\[
                 \mathcal F=F_*L
\]
has a perfect symmetric pairing
$\mathcal F\otimes\mathcal F\to L^{(2)}$ whose global pairing is
$\beta$.  This is finite Frobenius duality, locally the perfect
coefficient-of-$xy$ pairing in a two-element $2$-basis, with the
invariant two-form supplying the duality trivialization; see
\cite[Lemma~1.3.6, Proposition~1.3.7 and Theorem~1.3.8,
pp.~26--28]{BrionKumar2005}.

The bundle $\mathcal F$ is globally generated.  In the smooth-theta
case, Theorem~\ref{thm:ppav-three-theta-jets} gives second jets of
$3\Theta$ at every point.  A Frobenius fibre has local algebra
$k[x,y]/(x^2,y^2)$, which is a quotient of
$k[x,y]/(x,y)^3$, so these jets give precisely the required fibre
evaluation surjection.  In the product case the two degree-three
elliptic bundles evaluate surjectively on every length-two Frobenius
fibre, since the residual line bundle has degree one and has vanishing
$H^1$; tensor their evaluation maps.  These two cases exhaust the
principal surfaces by Lemma~\ref{lem:ppav-surface-dichotomy}.
Perfectness and global generation also show that $\beta$, and hence
the coefficient vector $\alpha$, is not zero.

Suppose $q_\alpha=0$.  The lemma supplies a nontrivial orthogonal
decomposition of $W$, and let $D$ project onto one three-dimensional
summand.  It preserves the kernel of every evaluation map to a fibre
of $\mathcal F$.  Indeed write a global section vanishing in that fibre
as $a+b$, with $a$ in that summand and $b$ in its orthogonal complement.
The value of $a$ pairs to zero with every section in its complement
by orthogonality.  It pairs to zero with every section in its own
summand because $a+b$ has zero value and $b$ is orthogonal to that
summand.  Global generation and perfectness therefore force the value
of $a$ to be zero, and then the value of $b$ is zero too.
The kernel of the surjective bundle evaluation
$W\otimes\OO_{A^{(2)}}\to\mathcal F$ is locally free.  The induced
map from that kernel to $\mathcal F$ vanishes in every closed fibre,
and hence vanishes identically.  Thus $D$ descends to an endomorphism
of $\mathcal F$, acting on global sections as the original rank-three
projector.

This contradicts Oda's simplicity theorem
\cite[Theorem~1.2(i), p.~45, with its preceding proof
pp.~43--45]{Oda1971}.  That theorem concerns isogenies of abelian
varieties of arbitrary dimension, including inseparable isogenies.
Its hypothesis here is satisfied because
$\lambda_L=3\lambda_\Theta$ restricts to an immersion on
$\ker F\subset A[2]$: $\lambda_\Theta$ is an isomorphism and three
is invertible on $A[2]$.  Therefore
$\End(F_*L)=k$.  A scalar endomorphism cannot induce that nontrivial
projector.  We have proved $q_\alpha\ne0$.

For clarity, the relative determinant, not just the full linearized
rank, is determined next.  The full coherent bundle on
$\mathbf P(W)$ is $W^*\otimes\OO(-1)$, and the relative subbundle
is $\Omega^1_{\mathbf P(W)}=\operatorname{ann}(s)$ under the Euler
sequence.  Its Frobenius-domain fibre is
$\operatorname{ann}(s^{[2]})$.  The contraction matrix is symmetric
and maps that domain into $\operatorname{ann}(s)$, since $M(s)s=0$.
On the affine chart $s_0=1$, use domain basis
$e_j+s_j^2e_0$ and target basis $e_i+s_ie_0$, $1\le i,j\le8$.
The resulting matrix has entries $M_{ij}+s_j^2M_{i0}$.
Expanding its determinant by these rank-one column additions and
using $\operatorname{adj}M=q_\alpha^2ss^t$ gives
\begin{equation}\label{eq:DE-A3-relative-theta-Hasse}
                 h_{\mathrm{rel}}(s)
                    =q_\alpha(s)^2\sum_{a\in I}s_a^3.
\end{equation}
Both factors are nonzero polynomials.  This is an actual determinant
on the rank-eight relative bundle, so its nonvanishing proves
bijectivity and stable rank eight.  Intersect this nonempty open with
the smooth-member open of the very ample system $|L|$.
Finally $g(Y)=1+L^2/2=10$, and exactness of stable parts in
\[
 0\longrightarrow H^1(A,\OO_A)\longrightarrow H^1(Y,\OO_Y)
   \longrightarrow E(Y,\OO_Y)\longrightarrow0
\]
gives $f_Y=f_A+8$.
\end{proof}

\begin{lemma}[Odd polarized norms preserve the curve defect]
\label{lem:DE-A3-ppav-odd-norm}
\status{Proved; root and monodromy full independent audits passed.}
For any principally polarized abelian surface $(A,\Theta)$ in
characteristic two and every odd $d\ge1$, the system $|3d\Theta|$
contains an integral nodal curve whose normalization is a smooth
degree-one theta-section of a principally polarized surface odd-isogenous
to $A$.  It can be chosen to have arithmetic-genus-minus-$2$-rank
defect exactly $2-f_A$.
\end{lemma}
\begin{proof}
Work geometrically.  Choose a maximal isotropic direct summand
in $A[d]\simeq(\mathbf Z/d\mathbf Z)^4$.  The quotient
$\psi:A\to A'$ has degree $d^2$, and $d\Theta$ descends to a
principal polarization $\Theta'$ on $A'$.
For the precise descent and principal-degree assertions see
\cite[Corollary~8.11, p.~117, and Proposition~8.19,
p.~119]{EdixhovenVanDerGeerMoonenAV}.
The polarized dual
\[
 \phi=\lambda_\Theta^{-1}\psi^\vee\lambda_{\Theta'}:A'\to A
\]
satisfies $\phi\psi=[d]$, $\psi\phi=[d]$,
$\deg\phi=d^2$ and $\phi^*\lambda_\Theta=d\lambda_{\Theta'}$.
It is finite \'etale since $d$ is odd.  Both surfaces have the
same $2$-rank $f_A$.

Let $L'=\OO_{A'}(3\Theta')$.  Proposition~\ref{prop:DE-A3-ppav-degree-one}
gives a nonempty open of smooth $C'\in|L'|$ with genus ten and
$2$-rank $8+f_A$.  We verify that a general such curve has nodal image
under $\phi$ without an ordinarity hypothesis on either surface.
A nonzero kernel translation preserving $L'$ acts nontrivially on
its projective space of sections, because $L'$ is very ample.
Its fixed locus is proper.  A translation not preserving $L'$ cannot
preserve a member.  Thus a general $C'$ has no nontrivial kernel
stabilizer, and $C'\to\phi(C')$ is generically one-to-one.

If $\Theta'$ is smooth, for a fixed nonzero kernel point $g$ the
rank-six failure of the two first jets at $P,P+g$ is contained in
\[
                   3P+g\in\Theta',\qquad
                   3P+2g\in\Theta'.
\]
The two theta translates are distinct, so their intersection is finite;
the map $[3]$ is finite.  Consequently only finitely many $P$ are
exceptional.  At every other $P$, two values and proportionality of
two nonzero gradients impose three independent conditions on sections,
while $P$ varies in dimension two.  At each exceptional point, the two
independent value conditions already give a proper locus.  This excludes
tangential double identifications for a general member.  The three-value
Lemma~\ref{lem:ppav-three-distinct-values} excludes triple
identifications: three distinct points $P,P+g,P+h$ impose three
conditions, with $P$ varying in dimension two.  There are finitely
many kernel elements and pairs.

For a product of elliptic curves the same assertions follow directly
from the cubic Segre embedding.  When both coordinates of $g$ are
nonzero, the crossed tangent exceptions are finite, by the cubic chord
law and finiteness of $[3]$.  When one coordinate is zero, the two first
jets have rank five.  Their only relation involves the two values and
the two derivatives in that factor; the transverse derivatives are
independent.  After setting both values to zero, at least one parallel
derivative can be nonzero: otherwise a nonzero section of a degree-three
elliptic bundle would vanish twice at two distinct points.  Proportionality
of the gradients is therefore one further independent equation.
Three distinct rank-one Segre tensors are collinear only if one factor
is constant and the other three points are collinear.  For three kernel
translates the latter condition is $3P_i+g_i+h_i=0$, a
one-dimensional locus on the surface.  Its value rank is two, and
elsewhere the rank is three.  The corresponding incidences are both
proper: their dimensions in the projective eight-space are at most
$8-2+1=7$ and $8-3+2=7$.  These arguments use only degree-three elliptic
bundles and the chord law, and so include supersingular factors.

Thus $Z=\phi(C')$ is integral nodal, with normalization $C'$.
For each nonzero kernel element the intersection
$C'\cap(C'-g)$ has degree $L'^2=18$.  No triple identifications
occur, so the number of nodes is $9(d^2-1)$.

The norm line bundle has Neron--Severi class $3d\Theta$:
push forward $\phi^*\Theta\equiv d\Theta'$ and use
$\deg\phi=d^2$ and torsion-freeness of the Neron--Severi group.
Choose symmetric representatives.  Both
$\operatorname{Nm}_\phi(L')$ and $\OO_A(3d\Theta)$ are symmetric,
so their quotient is a degree-zero two-torsion line bundle.
The homomorphism $3d\lambda_\Theta$ is an isomorphism on the geometric
two-torsion points.  Translation by one constant two-torsion point
therefore makes the image lie in the exact linear system $|3d\Theta|$.
This translation changes neither the curve nor its Frobenius cohomology.

The normalization sequence of an integral nodal curve adds a
Frobenius-bijective graph term of dimension its number of nodes.
It therefore preserves genus-minus-stable-rank defect.  The chosen
normalization has defect $10-(8+f_A)=2-f_A$, which proves the lemma.
For $d=1$ the isogeny is the identity and there are no nodes.
\end{proof}

\begin{lemma}[The connected Prym gives the even-degree bound]
\label{lem:DE-A3-ppav-even-bound}
\status{Proved; root and monodromy full independent audits passed.}
Let $Y\in|3d\Theta|$ be smooth on a principally polarized abelian
surface in characteristic two.  If $d$ is even, then
\[
                         g(Y)-f_Y\ge2(2-f_A).
\]
\end{lemma}
\begin{proof}
The restriction $j:\Pic^0(A)\to J_Y$ is a closed immersion.
For a nontrivial geometric degree-zero line bundle $P$ on $A$, one has
$H^0(P)=0$ and $H^1(P(-Y))=0$, the second by duality and ample
vanishing on an abelian variety.  The restriction sequence thus gives
$H^0(Y,P|_Y)=0$, so the geometric point kernel is trivial.
The differential of $j$ is injective because $H^1(A,\OO_A(-Y))=0$.
Its kernel is consequently a finite group scheme with zero tangent
space and no nontrivial geometric points, hence is trivial.

Identify $A$ with its dual through $\lambda_\Theta$, and let
$q:J_Y\to A$ be the polarized dual of this immersion.  It is
surjective and has smooth connected kernel $P_Y$.  One can see the
connectedness without assuming that an arbitrary isogeny has smooth
kernel: quotient $J_Y$ by the reduced identity component of $\ker q$.
The remaining map to $A$ is an isogeny whose dual kernel lies in
$\ker j=0$, so it has degree one by isogeny duality
\cite[Theorem~7.5, p.~100]{EdixhovenVanDerGeerMoonenAV}.

The induced-polarization formula is
\begin{equation}\label{eq:DE-A3-general-Prym-multiplier}
                       j^*\lambda_{J_Y}=(3d)\lambda_\Theta,
                       \qquad qj=[3d].
\end{equation}
This formula includes even $3d$; it is not a prime-to-characteristic
claim.  The exact norm/polarization calculation is
\cite[Section~3, equations~(3.2)--(3.3), pp.~346--347]{Debarre1994Degrees}:
for a curve numerically equivalent to $m\Theta$ on a principal surface,
its norm endomorphism is $[m]$ and its induced Jacobian polarization is
$m\Theta$.  Equivalently, restriction of the Jacobian pairing on
prime-to-two Tate modules is cup product with $[Y]=3d[\Theta]$;
faithfulness of that Tate realization identifies the homomorphisms in
\eqref{eq:DE-A3-general-Prym-multiplier}.  The polarized-dual convention
in this display avoids changing the sign when one identifies a Picard
variety with an Albanese variety.

If $d$ is even, $j$ restricts to an immersion of finite group schemes
$A[2]\hookrightarrow P_Y[2]$.  Twist this finite-group immersion by
Frobenius and apply functoriality of Verschiebung.  The identity
component of $\ker V_A$, of order $2^{2-f_A}$, embeds in the identity
component of $\ker V_{P_Y}$.  For any abelian variety $B$ the latter
order is $2^{\dim B-f_B}$, by the degree and etale-part formulas for
Verschiebung; see \cite[Propositions~5.20--5.22, pp.~80--81]
{EdixhovenVanDerGeerMoonenAV}.  Hence
\[
              \dim P_Y-f_{P_Y}\ge2-f_A.
\]
Finally $\dim J_Y=\dim P_Y+2$ and $f_Y=f_{P_Y}+f_A$.
The last equality follows also from Poincare reducibility and
isogeny invariance of $2$-rank.  Adding the two defects proves
$g(Y)-f_Y\ge2(2-f_A)$.
\end{proof}

\begin{theorem}[The universal parity formula for principal surfaces]
\label{thm:DE-A3-ppav-ranks}
\status{Proved; root and monodromy full independent audits passed.}
Let $(A,\Theta)$ be a principally polarized abelian surface over a
perfect field of characteristic two, with a representing line bundle,
and put $L=\OO_A(3\Theta)$.  Write $f_A\in\{0,1,2\}$ for its
geometric $2$-rank.  For every $d\ge1$, the geometric generic relative
stable rank in the complete system $|L^d|$ is
\begin{equation}\label{eq:DE-A3-ppav-rank-formula}
 \boxed{\qquad
 r_d=\begin{cases}
           9d^2-1,&d\text{ odd},\\
           9d^2-3+f_A,&d\text{ even}.
       \end{cases}\qquad}
\end{equation}
The relative coherent dimension is $9d^2-1$.  The full generic curve
defect is $2-f_A$ in odd degree and $2(2-f_A)$ in even degree.
Thus relative Frobenius is generically bijective exactly when $d$ is
odd or $A$ is ordinary.

For the fixed embedding by $L$ into $\mathbf P^8$, the same formula
holds for ambient degree-$d$ hypersurfaces whenever
\[
             d\ge\max\{1,\operatorname{reg}(\mathcal I_A)-1\}.
\]
For the cubic Segre embedding of a product of elliptic curves, it holds
for every $d\ge1$, with no restriction on their geometric
$j$-invariants.
\end{theorem}
\begin{proof}
Over an algebraic closure ample vanishing gives
$H^1(A,L^{-d})=0$, $h^2(A,L^{-d})=9d^2$ and $h^2(A,\OO_A)=1$.
The multiplication map on the last two groups is surjective, Serre-dual
to the inclusion of the defining nonzero section.  Its kernel is the
relative coherent quotient, of dimension $9d^2-1$.  This also constructs
a vector bundle of that rank on the whole coefficient projective space,
including the nodal boundary.  The full $H^1(\OO_Y)$ has constant
dimension $9d^2+1$.  Frobenius and all its iterates are bundle maps after
the appropriate Frobenius base changes, so lower bounds on stable rank
are open, by minors of an iterate at least as long as the bundle rank.

For odd $d$, Lemma~\ref{lem:DE-A3-ppav-odd-norm} supplies a nodal
member of full stable dimension
$9d^2+1-(2-f_A)=9d^2-1+f_A$.
Exactness of stable parts subtracts the ambient rank $f_A$, so the
relative Frobenius at that member is already bijective on all
$9d^2-1$ dimensions.  Its nonempty invertibility open meets the
nonempty smooth-member open.  This proves the odd formula.

For even $d\ge2$, choose generic sharp smooth odd-degree members
$C\in|L|$ and $D\in|L^{d-1}|$ meeting transversely.  They exist by
the odd case, and transversality is a nonempty open condition in the
two very ample systems.  Their connected nodal union has full defect
$2(2-f_A)$, since normalization adds only a Frobenius-bijective graph
term.  Its relative stable rank is therefore $9d^2-3+f_A$.
The open rank locus meets the smooth locus, giving that lower bound
generically.  Lemma~\ref{lem:DE-A3-ppav-even-bound} gives the matching
upper bound at every smooth member.  This proves
\eqref{eq:DE-A3-ppav-rank-formula} in every degree.

Regularity gives $H^1(\mathbf P^8,\mathcal I_A(d))=0$ in the displayed
range, so ambient sections surject onto the complete linear system.
The associated parameter map over the nonzero-section locus is an
affine bundle and has dense image, preserving the geometric generic
rank.  On a product, restriction of plane polynomials to either cubic
surjects in every positive degree, by the cubic hypersurface exact
sequence.  Their tensor products are restrictions of degree-$d$ Segre
polynomials, so the ambient map is surjective in every degree.
All the statements are geometric and are therefore preserved by
extending the original perfect field to its algebraic closure.
\end{proof}

\begin{remark}[Exact certificate and scope]
\label{rem:DE-A3-ppav-rank-checks}
The retained program
\path{agents/obstructions_crystals/code/A3_ppav_rank_checks.py}
checks the symbolic Pfaffian identity in all nine coordinates, the
forty coefficient points and their orthogonal-projector description,
the Fourier and quadratic-phase identities, and the exact isogeny,
node and parity dimension formulas.  The author run passed in
$0.27075901300122496$ seconds: nine symbolic Pfaffians expanded from
$945$ perfect matchings, all forty rank-three orthogonal projectors,
$1700$ supplementary relative-determinant values, and $30000$ dimension
identities.  Its complete output is retained alongside the program.
Root independently read the entire code and reran it in
$0.29324621799969464$ seconds, preserving the author output and saving
the separate result in \path{outputs/root_A3_ppav_rank_checks.json}.
The analytic proof above supplies the classification
over every algebraically closed characteristic-two field; a finite-field
sample is not used in place of that proof.

The earlier supersingular-product and generic-mixed formulas are
special cases.  In particular the algebraic ordinary-factor rank
refinement left explicit after Theorem~\ref{thm:DE-A3-mixed-generic}
is now covered by this theorem.
No structural rank formula for arbitrary smooth projective ambient
varieties is inferred, and this theorem alone does not determine the
full monodromy groups of nonordinary principal surfaces.
\end{remark}
